\documentclass[../../main.tex]{subfiles}
\graphicspath{{\subfix{../../image/}}}

\begin{document}

\section{其他函数空间}\label{section:ch05-其他函数空间}

\subsection{空间 $H^{-1}$}\label{subsection:ch05-H-1}

\begin{definition}\label{definition:ch05-H-1}
我们记 $H^{-1}(U)$ 为 $H^{1}_{0}(U)$ 的对偶空间。我们将写 $\langle,\rangle$ 表示 $H^{-1}(U)$ 与 $H^{1}_{0}(U)$ 之间的配对。
\end{definition}
\begin{remark}
换言之，$f$ 属于 $H^{-1}(U)$，只要 $f$ 是 $H^{1}_{0}(U)$ 上的有界线性泛函。请非常仔细地注意：我们并不把空间 $H^{1}_{0}$ 与其对偶等同。相反，正如我们马上将看到的，写
\begin{align*}
H^{1}_{0}(U)\subset L^{2}(U)\subset H^{-1}(U)
\end{align*}
更为有用。
\end{remark}

\begin{definition}\label{definition:ch05-H-1范数}
若 $f\in H^{-1}(U)$，定义其范数
\begin{align*}
\|f\|_{H^{-1}(U)}\triangleq\sup\left\{\langle f,u\rangle\mid u\in H^{1}_{0}(U),\ \|u\|_{H^{1}_{0}(U)}\leqslant 1\right\}.
\end{align*}
\end{definition}

\begin{theorem}[$H^{-1}$ 的刻画]\label{theorem:ch05-H-1刻画}
\begin{enumerate}[(i)]
\item 设 $f\in H^{-1}(U)$。则存在函数 $f_{0},f_{1},\cdots,f_{n}\in L^{2}(U)$ 使
\begin{align}\label{eq::jiwjio3398939erJFIOEJIE}
\langle f,v\rangle=\int_{U}f_{0}v+\sum_{i=1}^{n}f_{i}v_{x_{i}}\,\mathrm{d}x\qquad (v\in H^{1}_{0}(U)).
\end{align}

\item 此外，有
\begin{align*}
\|f\|_{H^{-1}(U)}=\inf\left\{\left(\int_{U}\sum_{i=0}^{n}|f_{i}|^{2}\,\mathrm{d}x\right)^{\frac{1}{2}}\ \middle|\ f\ \text{对}\ f_{0},\cdots,f_{n}\in L^{2}(U)\ \text{满足\eqref{eq::jiwjio3398939erJFIOEJIE}式}\right\}.
\end{align*}
\item 特别地，对一切 $u\in H^{1}_{0}(U)$，$f\in L^{2}(U)\subset H^{-1}(U)$，有
\begin{align}\label{eq::jioj382u82ur888238jjjj}
(f,u)_{L^{2}(U)}=\langle f,u\rangle.
\end{align}
这里$(,)_{L^{2}(U)}$表示$L^{2}(U)$中的内积.
\end{enumerate}
\end{theorem}
\begin{remark}
每当\eqref{eq::jiwjio3398939erJFIOEJIE}式成立时，我们记“$f=f_{0}-\sum_{i=1}^{n}f_{i,x_{i}}$”。
\end{remark}
\begin{proof}

1.
给定 $u,v\in H^{1}_{0}(U)$，定义内积
\begin{align*}
(u,v)_{H_0^1}\triangleq\int_{U}(Du\cdot Dv+uv)\,\mathrm{d}x.
\end{align*}
显然它是双线性的、对称的且正定的。该内积诱导的范数为
\begin{align*}
\|u\|_{H_0^1}\triangleq\left(\int_U(|Du|^2+|u|^2)\,\mathrm{d}x\right)^{\frac{1}{2}},
\end{align*}
这正是 $H_0^1(U)$ 的标准范数。由于 $H_0^1(U)$ 是完备的赋范线性空间，因此在这个内积下 $H_0^1(U)$ 是一个 Hilbert 空间。

设 $f\in H^{-1}(U)$，即 $f$ 是 $H_0^1(U)$ 上的连续线性泛函。对固定的 $f$，映射 $v\mapsto \langle f,v\rangle$ 是 $H_0^1(U)$ 上的连续线性泛函。由\hyperref[Functional Analysis-theorem:Riesz表示定理(Hilbert空间)]{Riesz表示定理}，存在唯一的 $u\in H_0^1(U)$，使得对一切 $v\in H_0^1(U)$，有
\begin{align}\label{eq:ch05-H-1-Riesz}
(u,v)_{H_0^1}=\langle f,v\rangle.
\end{align}
即
\begin{align}\label{eq:ch05-H-1-Riesz展开}
\int_{U}(Du\cdot Dv+uv)\,\mathrm{d}x=\langle f,v\rangle.
\end{align}
取 $f_0=u\in L^2(U)$，$f_i=u_{x_i}\in L^2(U)$ ($i=1,\cdots,n$)，则 \eqref{eq:ch05-H-1-Riesz展开} 式变为
\begin{align*}
\langle f,v\rangle=\int_{U}(f_0v+\sum_{i=1}^{n}f_iv_{x_i})\,\mathrm{d}x,
\end{align*}
这就证明了\eqref{eq::jiwjio3398939erJFIOEJIE}式。由于\hyperref[Functional Analysis-theorem:Riesz表示定理(Hilbert空间)]{Riesz表示定理}给出的 $u$ 是唯一的，故存在性得证。

2.
现设 $f\in H^{-1}(U)$ 有另一组表示
\begin{align}\label{eq:ch05-H-1-另一表示}
\langle f,v\rangle=\int_{U}(g_0v+\sum_{i=1}^{n}g_iv_{x_i})\,\mathrm{d}x
\end{align}
其中 $g_0,g_1,\cdots,g_n\in L^2(U)$。在 \eqref{eq:ch05-H-1-Riesz展开} 中令 $v=u$ 并利用 \eqref{eq:ch05-H-1-另一表示} 式，由Cauchy-Schwarz不等式推得
\begin{align*}
\int_{U}(|Du|^2+|u|^2)\,\mathrm{d}x
&\leqslant \left|\int_{U}(g_0u+\sum_{i=1}^{n}g_iu_{x_i})\,\mathrm{d}x\right|\\
&\leqslant \left(\int_{U}\sum_{i=0}^{n}|g_i|^2\,\mathrm{d}x\right)^{\frac{1}{2}}\left(\int_{U}(|u|^2+\sum_{i=1}^{n}|u_{x_i}|^2)\,\mathrm{d}x\right)^{\frac{1}{2}}\\
&=\left(\int_{U}\sum_{i=0}^{n}|g_i|^2\,\mathrm{d}x\right)^{\frac{1}{2}}\|u\|_{H_0^1}.
\end{align*}
两边同除以 $\|u\|_{H_0^1}$（若 $u=0$ 则不等式平凡成立），得到
\begin{align*}
\|u\|_{H_0^1}\leqslant\left(\int_{U}\sum_{i=0}^{n}|g_i|^2\,\mathrm{d}x\right)^{\frac{1}{2}}.
\end{align*}
注意到由\hyperref[Functional Analysis-theorem:Riesz表示定理(Hilbert空间)]{Riesz表示定理}得到的 $f_0=u$，$f_i=u_{x_i}$，故 $\|u\|_{H_0^1}^2=\int_{U}\sum_{i=0}^{n}|f_i|^2\,\mathrm{d}x$。因此式 \eqref{eq::jiwjio3398939erJFIOEJIE} 蕴含
\begin{align}\label{eq:ch05-H-1-极小性}
\left(\int_{U}\sum_{i=0}^{n}|f_i|^2\,\mathrm{d}x\right)^{\frac{1}{2}}\leqslant\left(\int_{U}\sum_{i=0}^{n}|g_i|^2\,\mathrm{d}x\right)^{\frac{1}{2}}.
\end{align}
这说明在所有的表示中，由\hyperref[Functional Analysis-theorem:Riesz表示定理(Hilbert空间)]{Riesz表示定理}给出的那组 $(f_0,\cdots,f_n)$ 使得 $\int \sum |f_i|^2$ 达到最小.

3.由 \eqref{eq::jiwjio3398939erJFIOEJIE} 得，当 $\|v\|_{H_0^1(U)}\leqslant 1$ 时，由Cauchy-Schwarz不等式得
\begin{align*}
|\langle f,v\rangle|
\leqslant \left(\int_{U}\sum_{i=0}^{n}|f_i|^2\,\mathrm{d}x\right)^{\frac{1}{2}}\left(\int_{U}(|v|^2+\sum_{i=1}^{n}|v_{x_i}|^2)\,\mathrm{d}x\right)^{\frac{1}{2}}\leqslant \left(\int_{U}\sum_{i=0}^{n}|f_i|^2\,\mathrm{d}x\right)^{\frac{1}{2}}.
\end{align*}
因此由算子范数的定义可得
\begin{align}\label{eq:ch05-H-1-范数上界}
\left\| f \right\| _{H^{-1}(U)}=\underset{\substack{
v\in H_{0}^{1}\left( U \right)\\
\left\| v \right\| _{H_{0}^{1}\left( U \right)}\leqslant 1\\
}}{\mathrm{sup}}\left| \left< f,v \right> \right|\leqslant \left( \int_U{\sum_{i=0}^n{|f_i|^2\,\mathrm{d}x}} \right) ^{\frac{1}{2}}.
\end{align}
另一方面，在 \eqref{eq:ch05-H-1-Riesz展开} 中令 $v=\frac{u}{\|u\|_{H_0^1}}$，推得
\begin{align*}
\left\langle f,\frac{u}{\|u\|_{H_0^1}}\right\rangle=\frac{1}{\|u\|_{H_0^1}}\int_U(|Du|^2+|u|^2)\,\mathrm{d}x=\|u\|_{H_0^1}.
\end{align*}
于是由算子范数的定义和$f_0=u,f_i=u_{x_i}$可得
\begin{align}\label{eq:ch05-H-1-范数下界}
\|f\|_{H^{-1}(U)}=\underset{\substack{
u\in H_{0}^{1}\left( U \right)\\
\left\| u \right\| _{H_{0}^{1}\left( U \right)}\leqslant 1\\
}}{\mathrm{sup}}\left| \left< f,u \right> \right|\geqslant\|u\|_{H_0^1}=\left(\int_{U}\sum_{i=0}^{n}|f_i|^2\,\mathrm{d}x\right)^{\frac{1}{2}}.
\end{align}
综合 \eqref{eq:ch05-H-1-范数上界} 与 \eqref{eq:ch05-H-1-范数下界}，得到
\begin{align*}
\|f\|_{H^{-1}(U)}=\left(\int_{U}\sum_{i=0}^{n}|f_i|^2\,\mathrm{d}x\right)^{\frac{1}{2}}.
\end{align*}
再结合 \eqref{eq:ch05-H-1-极小性} 可知，对所有满足\eqref{eq::jiwjio3398939erJFIOEJIE}式的 $f_0,\cdots,f_n$ 取下确界，就得到
\begin{align*}
\|f\|_{H^{-1}(U)}=\inf\left\{\left(\int_{U}\sum_{i=0}^{n}|f_i|^2\,\mathrm{d}x\right)^{\frac{1}{2}}\ \middle|\ f\ \text{对}\ f_0,\cdots,f_n\in L^2(U)\ \text{满足\eqref{eq::jiwjio3398939erJFIOEJIE}式}\right\}.
\end{align*}
断言 (ii) 得证。

4.首先 , 由 $H^{-1}(U)$ 的定义 , $H^{-1}(U)$ 是 $H_0^1(U)$ 的对偶空间 ,
因此若 $f\in H^{-1}(U)$ , 则 $f$ 是 $H_0^1(U)$ 上的连续线性泛函 , 其对偶配对记为 $\langle f,v\rangle$ , 其中 $v\in H_0^1(U)$ .

其次 , 说明 $L^2(U)$ 如何嵌入 $H^{-1}(U)$ . 对任意 $f\in L^2(U)$ , 定义映射 $T_f:H_0^1(U)\to\mathbb R$ 为
\begin{align*}
T_f(v)\triangleq\int_U f v\,\mathrm{d}x,\quad v\in H_0^1(U).
\end{align*}
由\hyperref[Real Analysis-theorem:实变函数--Hölder不等式]{Hölder不等式}得
\begin{align*}
|T_f(v)|
\leqslant \|f\|_{L^2(U)}\|v\|_{L^2(U)}
\leqslant \|f\|_{L^2(U)}\|v\|_{H_0^1(U)}.
\end{align*}
因此 $T_f$ 是 $H_0^1(U)$ 上的连续线性泛函 , 即 $T_f\in H^{-1}(U)$ . 这个嵌入 $f\mapsto T_f$ 就是 $L^2(U)\subseteq H^{-1}(U)$ 的含义 .

最后 , 在这个嵌入意义下 , 对任意 $u\in H_0^1(U)$ , 有
\begin{align*}
\langle f,u\rangle
=\langle T_f,u\rangle
=T_f(u)
=\int_U f u\,\mathrm{d}x.
\end{align*}
而 $L^2(U)$ 中的内积定义为
\begin{align*}
(f,u)_{L^2(U)}=\int_U f u\,\mathrm{d}x.
\end{align*}
因此得到
\begin{align*}
\langle f,u\rangle=\int_U f u\,\mathrm{d}x=(f,u)_{L^2(U)}.
\end{align*}
\end{proof}

\subsection{含时间的空间}\label{subsection:ch05-含时间空间}

下面我们研究另一些Sobolev空间，它们由把时间映入Banach空间的函数组成。这些空间在第 7 章构造线性抛物与双曲型 PDE 的弱解、以及第 9 章构造非线性抛物型 PDE 的弱解时将是本质的。

设 $X$ 为实Banach空间，范数为 $\|\,\cdot\,\|$。读者应先读\cref{section:appE-Banach值函数} 关于取值于 $X$ 的映射的测度与积分理论。

\begin{definition}\label{definition:ch05-Lp时间}
空间
\begin{align*}
L^{p}(0,T;X)
\end{align*}
由所有强可测函数 $u:[0,T]\to X$ 组成，满足
\begin{enumerate}[(i)]
\item 对 $1\leqslant p<\infty$，
\begin{align*}
\|u\|_{L^{p}(0,T;X)}\triangleq\left(\int_{0}^{T}\|u(t)\|^{p}\,\mathrm{d}t\right)^{1/p}<\infty;
\end{align*}
\item
\begin{align*}
\|u\|_{L^{\infty}(0,T;X)}\triangleq\operatorname*{ess\,sup}_{0\leqslant t\leqslant T}\|u(t)\|<\infty.
\end{align*}
\end{enumerate}
\end{definition}

\begin{definition}\label{definition:ch05-C时间}
空间
\begin{align*}
C([0,T];X)
\end{align*}
由所有连续函数 $u:[0,T]\to X$ 组成，满足
\begin{align*}
\|u\|_{C([0,T];X)}\triangleq\max_{0\leqslant t\leqslant T}\|u(t)\|<\infty.
\end{align*}
\end{definition}

\begin{definition}\label{definition:ch05-时间弱导数}
设 $u\in L^{1}(0,T;X)$。我们称 $v\in L^{1}(0,T;X)$ 为 $u$ 的弱导数，记作
\begin{align*}
u'=v,
\end{align*}
只要对一切实值检验函数 $\phi\in C^{\infty}_{c}(0,T)$ 有
\begin{align*}
\int_{0}^{T}\phi'(t)u(t)\,\mathrm{d}t=-\int_{0}^{T}\phi(t)v(t)\,\mathrm{d}t.
\end{align*}
\end{definition}

\begin{definition}\label{definition:ch05-W1p时间}
\leavevmode
\begin{enumerate}[(i)]
\item Sobolev空间
\begin{align*}
W^{1,p}(0,T;X)
\end{align*}
由所有函数 $u\in L^{p}(0,T;X)$ 组成，其中 $u'$ 在弱意义下存在且属于 $L^{p}(0,T;X)$。此外
\begin{align*}
\|u\|_{W^{1,p}(0,T;X)}\triangleq
\begin{cases}
\left(\int_{0}^{T}\|u(t)\|^{p}+\|u'(t)\|^{p}\,\mathrm{d}t\right)^{1/p}, & 1\leqslant p<\infty,\\
\operatorname*{ess\,sup}_{0\leqslant t\leqslant T}(\|u(t)\|+\|u'(t)\|), & p=\infty.
\end{cases}
\end{align*}
\item 我们写 $H^{1}(0,T;X)=W^{1,2}(0,T;X)$。
\end{enumerate}
\end{definition}

\begin{theorem}[抽象空间中的微积分]\label{theorem:ch05-抽象微积分}设 $u\in W^{1,p}(0,T;X)$，其中 $1\leqslant p\leqslant\infty$。则
\begin{enumerate}[(i)]
\item $u\in C([0,T];X)$（在一个零测集上重新定义之后）；
\item 对一切 $0\leqslant s\leqslant t\leqslant T$，
\begin{align*}
u(t)=u(s)+\int_{s}^{t}u'(\tau)\,\mathrm{d}\tau;
\end{align*}
\item 此外有估计
\begin{align}
\max_{0\leqslant t\leqslant T}\|u(t)\|\leqslant C\|u\|_{W^{1,p}(0,T;X)},
\end{align}
其中常数 $C$ 只依赖于 $T$。
\end{enumerate}
\end{theorem}
\begin{proof}

1. 把 $u$ 在 $(-\infty,0)$ 与 $(T,\infty)$ 上延拓为零，然后令 $u_{\varepsilon}=\eta_{\varepsilon}*u$，其中 $\eta_{\varepsilon}$ 为 $\mathbb{R}^{1}$ 上的通常磨光子。如\cref{theorem:ch05-光滑函数的局部逼近定理} 的证明中那样，我们检验在 $(\varepsilon,T-\varepsilon)$ 上 $(u_{\varepsilon})'=\eta_{\varepsilon}*u'$。

于是当 $\varepsilon\to 0$ 时，
\begin{align}
\begin{cases}
u_{\varepsilon}\to u\quad\text{于}\ L^{p}(0,T;X),\\
(u_{\varepsilon})'\to u'\quad\text{于}\ L^{p}_{\mathrm{loc}}(0,T;X).
\end{cases}
\end{align}
固定 $0<s<t<T$，计算
\begin{align*}
u_{\varepsilon}(t)=u_{\varepsilon}(s)+\int_{s}^{t}u_{\varepsilon}'(\tau)\,\mathrm{d}\tau.
\end{align*}
因而由上述收敛，对 a.e. $0<s<t<T$ 有
\begin{align}
u(t)=u(s)+\int_{s}^{t}u'(\tau)\,\mathrm{d}\tau.
\end{align}
由于映射 $t\mapsto\int_{0}^{t}u'(\tau)\,\mathrm{d}\tau$ 连续，断言 (i)、(ii) 随之成立。

2. 估计由上式容易推出。
\end{proof}

\begin{theorem}[更多微积分]\label{theorem:ch05-更多微积分}设 $u\in L^{2}(0,T;H^{1}_{0}(U))$，且 $u'\in L^{2}(0,T;H^{-1}(U))$。
\begin{enumerate}[(i)]
\item 则
\begin{align*}
u\in C([0,T];L^{2}(U))
\end{align*}
（在一个零测集上重新定义之后）。
\item 映射
\begin{align*}
t\mapsto\|u(t)\|_{L^{2}(U)}^{2}
\end{align*}
绝对连续，且对 a.e. $0\leqslant t\leqslant T$ 有
\begin{align*}
\frac{\mathrm{d}}{\mathrm{d}t}\|u(t)\|_{L^{2}(U)}^{2}=2\langle u'(t),u(t)\rangle.
\end{align*}
\item 此外有估计
\begin{align}
\max_{0\leqslant t\leqslant T}\|u(t)\|_{L^{2}(U)}
\leqslant C\left(\|u\|_{L^{2}(0,T;H^{1}_{0}(U))}+\|u'\|_{L^{2}(0,T;H^{-1}(U))}\right),
\end{align}
其中常数 $C$ 只依赖于 $T$。
\end{enumerate}
\end{theorem}
\begin{proof}

1. 先用反射再用截断函数，把 $u$ 延拓到更大的区间 $[-\sigma,T+\sigma]$（某个小 $\sigma>0$）；然后如前一证明那样定义正则化 $u_{\varepsilon}=\eta_{\varepsilon}*u$。于是对 $\varepsilon,\delta>0$，
\begin{align*}
\frac{\mathrm{d}}{\mathrm{d}t}\|u_{\varepsilon}(t)-u_{\delta}(t)\|_{L^{2}(U)}^{2}
=2\langle u_{\varepsilon}'(t)-u_{\delta}'(t),u_{\varepsilon}(t)-u_{\delta}(t)\rangle.
\end{align*}
因而
\begin{align}
\|u_{\varepsilon}(t)-u_{\delta}(t)\|_{L^{2}(U)}^{2}
=\|u_{\varepsilon}(s)-u_{\delta}(s)\|_{L^{2}(U)}^{2}
+2\int_{s}^{t}\langle u_{\varepsilon}'(\tau)-u_{\delta}'(\tau),u_{\varepsilon}(\tau)-u_{\delta}(\tau)\rangle\,\mathrm{d}\tau
\end{align}
对一切 $0\leqslant s,t\leqslant T$。固定任一点 $s\in(0,T)$ 使 $u_{\varepsilon}(s)\to u(s)$ 于 $L^{2}(U)$。于是上式蕴含
\begin{align*}
\limsup_{\varepsilon,\delta\to 0}\sup_{0\leqslant t\leqslant T}\|u_{\varepsilon}(t)-u_{\delta}(t)\|_{L^{2}(U)}^{2}
\leqslant\lim_{\varepsilon,\delta\to 0}\int_{0}^{T}\|u_{\varepsilon}'(\tau)-u_{\delta}'(\tau)\|_{H^{-1}(U)}^{2}
+\|u_{\varepsilon}(\tau)-u_{\delta}(\tau)\|_{H^{1}_{0}(U)}^{2}\,\mathrm{d}\tau=0.
\end{align*}
因此磨光函数族 $\{u_{\varepsilon}\}_{0<\varepsilon\leqslant 1}$ 在 $C([0,T];L^{2}(U))$ 中收敛于某极限 $v\in C([0,T];L^{2}(U))$。由于我们还知道对 a.e. $t$ 有 $u_{\varepsilon}(t)\to u(t)$，推得 $u=v$ a.e.

2. 我们类似地有
\begin{align*}
\|u_{\varepsilon}(t)\|_{L^{2}(U)}^{2}
=\|u_{\varepsilon}(s)\|_{L^{2}(U)}^{2}
+2\int_{s}^{t}\langle u_{\varepsilon}'(\tau),u_{\varepsilon}(\tau)\rangle\,\mathrm{d}\tau,
\end{align*}
于是把 $u$ 与上面的 $v$ 等同，得
\begin{align}
\|u(t)\|_{L^{2}(U)}^{2}
=\|u(s)\|_{L^{2}(U)}^{2}
+2\int_{s}^{t}\langle u'(\tau),u(\tau)\rangle\,\mathrm{d}\tau
\end{align}
对一切 $0\leqslant s,t\leqslant T$。

3. 为得到估计，我们对上式关于 $s$ 积分，回忆不等式 $|\langle u',u\rangle|\leqslant\|u'\|_{H^{-1}(U)}\|u\|_{H^{1}_{0}(U)}$，并作一些简单估计。
\end{proof}

\begin{theorem}[映入更好空间的映射]\label{theorem:ch05-更好空间映射}设 $U$ 开、有界，$\partial U$ 光滑。取 $m$ 为非负整数。设 $u\in L^{2}(0,T;H^{m+2}(U))$，且 $u'\in L^{2}(0,T;H^{m}(U))$。
\begin{enumerate}[(i)]
\item 则
\begin{align*}
u\in C([0,T];H^{m+1}(U))
\end{align*}
（在一个零测集上重新定义之后）。
\item 此外有估计
\begin{align}
\max_{0\leqslant t\leqslant T}\|u(t)\|_{H^{m+1}(U)}
\leqslant C\left(\|u\|_{L^{2}(0,T;H^{m+2}(U))}+\|u'\|_{L^{2}(0,T;H^{m}(U))}\right),
\end{align}
其中常数 $C$ 只依赖于 $T$、$U$ 与 $m$。
\end{enumerate}
\end{theorem}
\begin{proof}

1. 先设 $m=0$，此时
\begin{align*}
u\in L^{2}(0,T;H^{2}(U)),\qquad u'\in L^{2}(0,T;L^{2}(U)).
\end{align*}
我们选有界开集 $V\Supset U$，然后如\cref{theorem:ch05-扩张定理} 构造相应延拓 $\overline{u}=Eu$。鉴于该节的估计，我们有 $\overline{u}\in L^{2}(0,T;H^{2}(V))$，且
\begin{align}
\|\overline{u}\|_{L^{2}(0,T;H^{2}(V))}\leqslant C\|u\|_{L^{2}(0,T;H^{2}(U))},
\end{align}
其中 $C$ 为适当常数。此外 $\overline{u}'\in L^{2}(0,T;L^{2}(V))$，并有估计
\begin{align}
\|\overline{u}'\|_{L^{2}(0,T;L^{2}(V))}\leqslant C\|u'\|_{L^{2}(0,T;L^{2}(U))}.
\end{align}
这只要考虑 $t$ 变量中的差商，回忆\cref{subsection:ch05-差商} 中的方法，并注意 $E$ 是从 $L^{2}(U)$ 到 $L^{2}(V)$ 的有界线性算子，即得。

2. 暂时设 $\overline{u}$ 光滑。我们计算
\begin{align*}
\left|\frac{\mathrm{d}}{\mathrm{d}t}\left(\int_{V}|D\overline{u}|^{2}\,\mathrm{d}x\right)\right|
=2\left|\int_{V}D\overline{u}\cdot D\overline{u}'\,\mathrm{d}x\right|
=2\left|\int_{V}\Delta\overline{u}\,\overline{u}'\,\mathrm{d}x\right|
\leqslant C\left(\|\overline{u}\|_{H^{2}(V)}^{2}+\|\overline{u}'\|_{L^{2}(V)}^{2}\right).
\end{align*}
分部积分时没有边界项，因为延拓 $\overline{u}=Eu$ 在 $V$ 中有紧支集。积分并回忆上述两估计，得
\begin{align}
\max_{0\leqslant t\leqslant T}\|u(t)\|_{H^{1}(U)}
\leqslant C\left(\|u\|_{L^{2}(0,T;H^{2}(U))}+\|u'\|_{L^{2}(0,T;L^{2}(U))}\right).
\end{align}
若 $u$ 不光滑，如前用 $u_{\varepsilon}\triangleq\eta_{\varepsilon}*u$ 逼近，我们得到同一估计。如前一证明中那样，还推得 $u\in C([0,T];H^{1}(U))$。

3. 在 $m\geqslant 1$ 的一般情形，取 $\alpha$ 为阶 $|\alpha|\leqslant m$ 的多重指标，令 $v\triangleq D^{\alpha}u$。则
\begin{align*}
v\in L^{2}(0,T;H^{2}(U)),\qquad v'\in L^{2}(0,T;L^{2}(U)).
\end{align*}
我们把估计~\eqref{align:ch05-H1估计} 应用于 $v$（以 $v$ 代 $u$），并对所有 $|\alpha|\leqslant m$ 求和，即导出估计~\eqref{align:ch05-H1估计}。
\end{proof}

\end{document}
